\uuid{sTqd}
\titre{Ensembles}
\auteur{bodin ; croizat ; sacre}
\organisation{exo7}
\licence{CC BY-NC-SA 4.0}
\datecreate{2022-03-08}
\niveau{L1}
\module{Algèbre}
\chapitre{Logique, ensemble, raisonnement}
\sousChapitre{Ensemble}
\difficulte{1}
\theme{Ensembles}
\exo7qcmid{hilisit.logique.020}

\contenu{
\begin{question}
\qclassification{100.02}
Soit la fonction $f : \Rr \to \Rr$ définie par $f(x) = x^2+2$.
Quelles sont les affirmations vraies ?
\begin{answers}
	\good{L'image de $-2$ est $6$.}
	\bad{Un antécédent de $18$ est $-5$.}
	\bad{La valeur $2$ admet plusieurs images.}
	\good{La valeur $18$ admet plusieurs antécédents.}
\end{answers}
\begin{explanations}
L'image de $-2$ est $6$ car $f(-2) = 6$. \\
Le nombre $-5$ n'est pas antécédent de $18$ car $f(-5) = 27 \neq 18$. \\
La valeur $x=2$ admet une seule image. C'est en fait vrai pour tout $x\in\Rr$ !\\
Les antécédents de $18$ sont les solutions de l'équations $f(x) = 18$, c'est-à-dire $x^2+2=18$ ou encore $x^2=16$, ce sont donc $x=4$ et $x=-4$. Il y a donc deux antécédents à la valeur $18$.
\end{explanations}
\end{question}
}
